\documentclass[10pt,twocolumn]{article}
\usepackage[T1]{fontenc}
\usepackage[utf8]{inputenc}
\usepackage{lmodern}
\usepackage[margin=1.8cm,top=2.4cm,bottom=2cm,columnsep=0.6cm]{geometry}
\usepackage{fancyhdr}
\usepackage{titlesec}
\usepackage{booktabs}
\usepackage{amsmath,amssymb,amsfonts}
\usepackage{microtype}
\usepackage{xcolor}
\usepackage{mdframed}
\usepackage{parskip}
\usepackage{enumitem}
\usepackage{hyperref}

\definecolor{rulered}{RGB}{180,30,30}
\definecolor{boxgray}{RGB}{245,245,245}
\definecolor{keywordgray}{RGB}{100,100,100}

\pagestyle{fancy}
\fancyhf{}
\fancyhead[L]{\textsc{\small Claude Mag}}
\fancyhead[C]{\small\textit{Machine Learning Research Digest}}
\fancyhead[R]{\small 2026-09-23}
\fancyfoot[C]{\small\thepage}
\renewcommand{\headrulewidth}{0.8pt}

\titleformat{\section}{\normalsize\bfseries\color{rulered}}{}{0pt}{}%
  [\vspace{-4pt}\color{rulered}\rule{\columnwidth}{0.4pt}]
\titlespacing{\section}{0pt}{8pt}{4pt}

\hypersetup{colorlinks=true,urlcolor=rulered,linkcolor=black}

\begin{document}
\setlength{\parindent}{0pt}
\setlength{\parskip}{4pt}

{\small\color{keywordgray}\textbf{Keywords:}
  LLM reasoning \textbullet{}
  formal verification \textbullet{}
  SMT solver \textbullet{}
  chain-of-thought}\par\vspace{4pt}

{\LARGE\bfseries\raggedright
  LogicTrack}\par\vspace{4pt}

{\large\itshape\raggedright
  Auditing LLM Reasoning Trajectories with Formal Logic Solvers and
  Step-Level Backtracking Rewards}\par\vspace{6pt}

{\small\textbf{By} Jingyu Hu, Shu Yang, Weiru Liu \& Di Wang}\par\vspace{2pt}

{\small\color{keywordgray}arXiv:2609.21492 \textbullet{} September 2026}
\par\vspace{2pt}\noindent{\color{rulered}\rule{\columnwidth}{0.6pt}}\vspace{6pt}

LogicTrack auto-formalizes each chain-of-thought step into an SMT-LIB
specification, verifies it with a Z3-class solver, and converts the
diagnostics into a Solver-Based Backtracking Reward that guides both
inference-time tree search and supervised fine-tuning---improving
step-level verifiability across 8 benchmarks and 7 LLMs.

\begin{mdframed}[backgroundcolor=boxgray,linecolor=rulered,linewidth=1pt,
    innerleftmargin=6pt,innerrightmargin=6pt,
    innertopmargin=6pt,innerbottommargin=6pt]
\textbf{\small AT A GLANCE}\par\vspace{4pt}
\begin{description}[leftmargin=0pt,labelsep=4pt,itemsep=2pt,parsep=0pt,
    font=\small\bfseries]
  \item[Problem:] RLVR optimises final-answer correctness but leaves the
    logical validity of intermediate reasoning steps unverified and
    potentially unsound.
  \item[Why it matters:] Models that reach correct answers via logically
    invalid steps are brittle, untrustworthy, and opaque to auditors.
  \item[Method:] Decompose each CoT step into context, assumptions, and
    conclusion; auto-formalize to SMT-LIB; verify with a theorem prover;
    combine into a Solver-Based Backtracking Reward (SBR).
  \item[Result:] Step-level verifiability improves 15--30 pp; final
    accuracy preserved or improved across 8 benchmarks, 7 LLMs.
  \item[Evidence:] Systematic evaluation on GSM8K, MATH, ARC, LogiQA,
    FOLIO, ProntoQA, and custom multi-step deduction suites.
\end{description}
\end{mdframed}

\section{The Big Picture}

Large language model reasoning has been transformed by chain-of-thought
(CoT) prompting and reinforcement learning with verifiable rewards (RLVR).
Models now generate extended thinking traces before answering, and these
traces substantially improve performance on complex tasks. But the training
signal in RLVR is a binary or scalar reward over the final answer---not
over the intermediate steps.

This creates a subtle failure mode: models can learn to produce long,
plausible-looking reasoning chains that are internally inconsistent or
logically invalid, while still arriving at correct answers through
pattern-matching or lucky leaps. Such traces are unauditable, brittle to
distributional shift, and dangerous in high-stakes domains where the
reasoning process---not just the answer---must be verified.

LogicTrack addresses this gap with a neuro-symbolic framework that
provides free, scalable, per-step logical verification as a training and
inference-time signal.

\section{Why Existing Approaches Fall Short}

\textbf{Outcome-based RLVR} (GRPO, PPO-R1) provides a single reward per
trajectory, giving no credit for valid intermediate steps or penalty for
invalid ones. A model is rewarded equally for a logically sound path and
a lucky leap.

\textbf{Process reward models (PRMs)} score intermediate steps but are
trained discriminatively on expensive human labels and do not carry
logical soundness guarantees---they reflect human preferences, not logical
validity.

\textbf{Self-consistency and majority voting} increase answer reliability
through repeated sampling but do not verify any individual trace.

\textbf{Direct SMT encoding} cannot process natural language without a
formalization layer, which has historically required domain-specific
engineering and failed to generalize.

LogicTrack pairs a neural auto-formalization module (an LLM translator)
with a classical SMT solver, getting the generalization of neural models
and the soundness guarantees of symbolic reasoning.

\section{Core Mechanism}

\textbf{Step Decomposition.} Each natural-language step $s_k$ is
decomposed into three fields:
\begin{enumerate}[itemsep=2pt,parsep=0pt]
  \item \textbf{Context}: facts established by prior steps $s_1,\ldots,s_{k-1}$
  \item \textbf{Background assumptions}: domain knowledge or definitions
  \item \textbf{Conclusion}: the claimed logical consequence of 1 and 2
\end{enumerate}

\textbf{Auto-formalization.} An LLM translator maps each triple into
an SMT-LIB specification covering linear arithmetic, equality, and
propositional logic---the fragment that covers the large majority of
mathematical and logical reasoning steps.

\textbf{Solver Verification.} A Z3-class SMT solver checks the
specification and returns one of:
\textit{syntax error}, \textit{satisfiable}, \textit{entailed}
(conclusion follows from context + assumptions),
\textit{refuted} (conclusion contradicts them), or
\textit{undecided} (solver timeout).

\textbf{Solver-Based Backtracking Reward (SBR).} SBR combines:
\begin{itemize}[itemsep=2pt,parsep=0pt]
  \item A \emph{premise-fidelity score} measuring how faithfully the
    SMT context matches prior step conclusions
  \item A \emph{solver verification score}: entailed $>$ satisfiable
    $>$ undecided $>$ refuted
\end{itemize}

\textbf{Inference-time Backtracking.} SBR guides tree search: when the
current path stalls, the search backtracks to the step with the
lowest SBR score and regenerates from that branch.

\textbf{SFT data construction.} Successful backtracking trajectories
are formatted as SFT examples, teaching fine-tuned models to internalise
step-wise auditing without requiring the solver at inference time.

\section{Technical Formulation}

Let trajectory $\tau = (s_1,\ldots,s_T, a)$ where $a$ is the final
answer. Each step $s_k$ maps to an SMT triple
$(\mathrm{ctx}_k, \mathrm{bg}_k, \mathrm{conc}_k)$.

The SMT validity check:
\[
\mathrm{SMT\text{-}check}: \;\mathrm{ctx}_k \wedge \mathrm{bg}_k
  \;\models\; \mathrm{conc}_k
\]

Step-level SBR:
\[
\mathrm{SBR}(s_k)
= \alpha\cdot\mathrm{Fid}(s_k)
+ (1-\alpha)\cdot\mathrm{SolverScore}(s_k)
\]

Trajectory-level SBR (used as training reward):
\[
\mathrm{SBR}(\tau)
= \frac{1}{T}\sum_{k=1}^{T}\mathrm{SBR}(s_k)
  \cdot \mathbf{1}[a = a^*]
\]

Backtrack point selection:
\[
k^* = \operatorname*{arg\,min}_{k}\,\mathrm{SBR}(s_k)
\]
The search then regenerates $s_{k^*},\ldots,s_T$ from the new branch.

\section{Learning or Inference Procedure}

\textbf{Inference-time:}
\begin{enumerate}[itemsep=2pt,parsep=0pt]
  \item Generate initial trajectory $\tau$.
  \item Score each step with SBR; if $\mathrm{SBR}(\tau)<\delta$
    backtrack to $k^*$ and regenerate.
  \item Repeat until answer correct or budget exhausted.
\end{enumerate}

\textbf{Training-time SFT:}
\begin{enumerate}[itemsep=2pt,parsep=0pt]
  \item Collect successful backtracking trajectories.
  \item Format with explicit \texttt{<backtrack>} markers.
  \item Fine-tune base model; solver not required at inference.
\end{enumerate}

\section{What the Guarantee Says}

For steps in the linear-arithmetic and propositional-logic fragment,
the SMT check is sound and (conditionally) complete: an \emph{entailed}
verdict guarantees the conclusion logically follows from the stated
context and assumptions. Steps marked \emph{refuted} are provably
invalid under the formalized premises. The \emph{undecided} category
handles solver incompleteness gracefully.

\section{Experimental Findings}

\begin{itemize}[itemsep=2pt,parsep=0pt]
  \item \textbf{8 benchmarks:} GSM8K, MATH, ARC-Challenge, LogiQA,
    FOLIO, ProntoQA, HellaSwag-hard, custom multi-step deduction
  \item \textbf{7 LLMs:} 7B to 70B parameter range
  \item Step-level verifiability improves by 15--30 percentage points
    across all model families
  \item Final answer accuracy improves on 5 of 8 benchmarks;
    negligible change on the remaining 3
  \item SFT on backtracking traces provides lasting improvements without
    runtime solver dependency
  \item Largest gains on tasks requiring deductive chaining (FOLIO, ProntoQA)
\end{itemize}

\section{Ablations and Interpretation}

\textbf{SBR vs.\ outcome-only reward:} Outcome reward improves final
accuracy but not step-level verifiability; SBR improves both, confirming
that process supervision is strictly more informative.

\textbf{Backtracking vs.\ repeated sampling:} At equal inference budget,
backtracking achieves higher accuracy by directing search toward logical
faults rather than sampling independently.

\textbf{Auto-formalization quality:} $\approx$85\% of steps are
successfully formalized; failures are handled as \emph{undecided} and
penalised only lightly, preventing over-penalisation of complex steps
that exceed the SMT-LIB fragment.

\textbf{$\alpha$ sensitivity:} Performance is stable for
$\alpha \in [0.3, 0.7]$, with a slight preference for $\alpha = 0.5$.

\section*{Reference}

Jingyu Hu, Shu Yang, Weiru Liu, Di Wang.
\textbf{LogicTrack: Auditing Reasoning Trajectories of Large Language
Models with Formal Logic Solvers}.
arXiv:2609.21492, September 2026.\\
\url{https://arxiv.org/abs/2609.21492}

\end{document}
